\section{Revisi\'on bibliogr\'afica}

\subsection{Modelos fundamentales de optimizaci\'on de portafolios}

La optimizaci\'on de portafolios tiene su origen en el trabajo seminal de
Markowitz \cite{markowitz1952}, que formaliz\'o el principio de
diversificaci\'on mediante el modelo media-varianza. Su propuesta establece
que un inversionista racional selecciona portafolios sobre la frontera
eficiente, es decir, aquellos que maximizan el retorno esperado para un nivel
de riesgo dado. La principal limitaci\'on del modelo es su sensibilidad a las
estimaciones de $\mu$ y $\Sigma$: peque\~nos errores en los retornos
esperados pueden generar portafolios concentrados o irrazonables
\cite{michaud1989}. En el presente proyecto, esta fragilidad se observa
emp\'iricamente: el portafolio Ideal de mercado de Markowitz exhibe un Sharpe
te\'orico de 5,01 que colapsa a 1,80 fuera de muestra con datos de pandemia,
confirmando el fen\'omeno de \textit{overfitting} documentado por Michaud.

El modelo CAPM de Sharpe \cite{sharpe1964} extiende el trabajo de Markowitz
introduciendo un modelo de equilibrio de mercado, donde el retorno esperado
de un activo depende \'unicamente de su exposici\'on al riesgo sistem\'atico,
medido por el coeficiente $\beta$. El modelo multifactorial de Fama y French
\cite{famafrench1993} incorpora dos factores adicionales, tama\~no (SMB) y
valor (HML), para capturar anomal\'ias emp\'iricas. Aunque estos modelos no
se implementan directamente en FinPUC, informan la construcci\'on del prior
de equilibrio en Black-Litterman: los retornos impl\'icitos $\pi = \delta
\Sigma w_{\text{market}}$ asumen que el mercado est\'a en equilibrio, lo que
es consistente con la l\'ogica del CAPM.

\subsection{El modelo Black-Litterman}

Black y Litterman \cite{black1992} propusieron un enfoque bayesiano que
combina el equilibrio impl\'icito del mercado con expectativas subjetivas del
inversionista. El prior se obtiene invirtiendo el problema de Markowitz: dado
que el mercado est\'a en equilibrio, los retornos impl\'icitos son
$\pi = \delta \Sigma w_{\text{market}}$. Las \textit{views} se incorporan
mediante una actualizaci\'on bayesiana que produce estimaciones m\'as estables
y portafolios menos extremos. He y Litterman \cite{he1999} documentaron la
implementaci\'on pr\'actica del modelo, destacando su capacidad para generar
portafolios intuitivamente razonables sin restricciones artificiales de peso.
Idzorek \cite{idzorek2005} extendi\'o el marco proponiendo m\'etodos para
calibrar la confianza de cada \textit{view}, problema que en FinPUC se aborda
fijando $c = 0.50$ para las \textit{views} de momentum y $c = 0.35$ para la
de desempleo.

\subsection{Estimaci\'on de covarianzas y medidas de riesgo}

Ledoit y Wolf \cite{ledoit2004} propusieron el m\'etodo de \textit{shrinkage}
para la estimaci\'on de matrices de covarianza de grandes dimensiones,
demostrando que la combinaci\'on convexa entre la covarianza muestral y un
estimador estructurado reduce el error de estimaci\'on fuera de muestra. En
FinPUC se implementa con $\lambda = 0.20$, valor seleccionado para
estabilizar una matriz de $601 \times 601$ sin perder demasiada estructura de
correlaci\'on. La elecci\'on de $\lambda$ no fue optimizada formalmente y
constituye un supuesto del modelo.

Rockafellar y Uryasev \cite{rockafellar2000} introdujeron el Valor en Riesgo
Condicional (CVaR) como medida coherente de riesgo, superando las limitaciones
del VaR est\'andar al cuantificar la severidad promedio de las p\'erdidas en
la cola. En el proyecto, el CVaR al 95\% se utiliza como indicador de
desempe\~no complementario al Sharpe, no como restricci\'on del optimizador.
Esta decisi\'on responde a que la implementaci\'on con CVXPY utiliza cotas de
volatilidad por perfil, que son computacionalmente m\'as eficientes para la
restricci\'on de 10 segundos.

\subsection{Sistemas de asesor\'ia financiera automatizada}

En el contexto de plataformas de \textit{robo-advisory}, D'Acunto, Prabhala y
Rossi \cite{dacunto2019} documentaron c\'omo la automatizaci\'on ha
democratizado el acceso a estrategias de inversi\'on, permitiendo
recomendaciones personalizadas a bajo costo operacional. FinPUC se enmarca en
esta l\'inea, con la particularidad de que incorpora expl\'icitamente el
comportamiento del cliente (funciones $P_1$ y $P_2$) y la utilidad de la
empresa (comisiones $k\%$) como variables del modelo, lo que va m\'as all\'a
de la optimizaci\'on financiera pura.

\subsection{Fortalezas y debilidades de la metodolog\'ia seleccionada}

La combinaci\'on de Black-Litterman con Monte Carlo resulta adecuada para el problema de FinPUC porque une dos necesidades del sistema recomendador: construir portafolios m\'as estables que Markowitz puro y evaluar su impacto en variables de decisi\'on para cliente y empresa. Black-Litterman reduce la dependencia de estimaciones hist\'oricas extremas al anclar los retornos esperados al equilibrio de mercado. Esto se observa en los resultados: con datos de pandemia, la mejor variante BL alcanza Sharpe 2,38, mientras que Markowitz cae a 2,02.

Otra fortaleza es que el modelo permite incorporar \textit{views} de forma controlada. En vez de reemplazar completamente el prior de mercado, las se\~nales de momentum, desempleo o capitalizaci\'on ajustan los retornos esperados seg\'un su nivel de confianza. Luego, Monte Carlo traduce esos portafolios a escenarios desfavorable, neutro y favorable, permitiendo medir riqueza terminal, retiro del cliente y utilidad para FinPUC sin depender solo de una m\'etrica financiera.

La principal debilidad es que la calidad del resultado depende directamente de la calidad de las \textit{views}. En particular, la view de desempleo se modela con un supuesto fijo y no con una serie macro observada, lo que limita su validez predictiva. Adem\'as, par\'ametros como $\tau = 0.05$ y la matriz $\Omega = P\tau\Sigma P^\top / c$ requieren sensibilidad adicional, porque cambios en la confianza asignada a las views pueden modificar el ranking final.

Finalmente, la simulaci\'on Monte Carlo asume retornos normales, aunque el an\'alisis distribucional rechaza normalidad en todos los activos. Esto puede subestimar eventos extremos y generar tasas de retiro demasiado optimistas. Tambi\'en queda pendiente incorporar costos de transacci\'on en la optimizaci\'on, ya que actualmente aparecen solo como comisiones en P4. Sin esa penalizaci\'on, algunas carteras podr\'ian verse atractivas en el modelo pero ser menos eficientes en una implementaci\'on real con turnover alto.


% ============================================================================
% 4.1 Arquitectura general
% ============================================================================
\section{Arquitectura general del sistema}

El sistema se estructura en tres capas metodol\'ogicas de complejidad
creciente, ejecutadas secuencialmente. La primera capa implementa la
optimizaci\'on media-varianza de Markowitz como l\'inea base anal\'itica. La
segunda extiende Markowitz con el modelo Black-Litterman, que incorpora un
prior de equilibrio de mercado y \textit{views} cuantitativas. La tercera capa
es una simulaci\'on Monte Carlo que traduce las m\'etricas financieras de las
capas anteriores en proyecciones de riqueza del cliente y utilidad de la
empresa. Adicionalmente, un benchmark equiponderado funciona como caso base
independiente de par\'ametros estimados. En total, el sistema compara seis
t\'ecnicas: Equiponderado, Markowitz, y cuatro variantes de Black-Litterman
(momentum general, momentum por capitalizaci\'on a 6 meses, momentum por
capitalizaci\'on a 1 a\~no, y desempleo macro).

% ============================================================================
% 4.2 Estimación de parámetros
% ============================================================================
\section{Estimaci\'on de par\'ametros}

\subsection{Retornos esperados}

Los retornos esperados $\mu \in \mathbb{R}^n$ se estiman como promedios
diarios hist\'oricos de los retornos totales de cada acci\'on en la ventana de
entrenamiento.

\subsection{Matriz de covarianzas con \textit{shrinkage}}

La matriz de covarianzas muestral $\Sigma_{\text{muestral}}$ se estabiliza
mediante \textit{shrinkage} lineal hacia su diagonal
\cite{ledoit2004}:
\begin{equation}
\Sigma_{\text{shrink}} = (1 - \lambda)\,\Sigma_{\text{muestral}}
+ \lambda \cdot \text{diag}(\Sigma_{\text{muestral}}),
\end{equation}
con $\lambda = 0.20$. Esto reduce el error de estimaci\'on fuera de muestra en
matrices de alta dimensi\'on ($n = 601$), donde la covarianza muestral tiende
a sobreajustar correlaciones espurias. La matriz resultante se fuerza a ser
semidefinida positiva mediante proyecci\'on espectral
(\texttt{nearest\_psd}), truncando valores propios negativos a cero. Este paso
es requisito para que los problemas cuadr\'aticos resueltos por CVXPY sean
convexos.

\subsection{Anualizaci\'on}

Los retornos se anualizan asumiendo 252 d\'ias h\'abiles bursátil por a\~no:
$R_{\text{anual}} = \bar{r}_{\text{diario}} \times 252$. La volatilidad se
escala como $\sigma_{\text{anual}} = \sigma_{\text{diaria}} \times
\sqrt{252}$. La tasa libre de riesgo se fija en $R_f = 2\%$ anual.

% ============================================================================
% 4.3 Optimización media-varianza de Markowitz
% ============================================================================
\section{Optimizaci\'on media-varianza de Markowitz}

\subsection{Formulaci\'on por perfil de riesgo}

Para cada perfil de riesgo con par\'ametro de aversi\'on $\gamma$ y cota de
volatilidad $\sigma_{\max}^2$, se resuelve el problema cuadr\'atico:
\begin{equation}
\max_{w} \quad \mu^\top w - \frac{1}{2}\,\gamma\, w^\top \Sigma\, w
\end{equation}
sujeto a:
\begin{align}
\sum_{i=1}^{n} w_i &= 1 \\
0 \leq w_i &\leq 0.20 \quad \forall\, i = 1, \ldots, n \\
w^\top \Sigma\, w &\leq \sigma_{\max}^2
\quad \text{(cuando el perfil lo requiere)}
\end{align}

La restricci\'on $w_i \geq 0$ impide ventas cortas; la suma igual a 1
invierte todo el capital; y la cota superior $w_{\max} = 0.20$ evita
concentraci\'on extrema en un solo activo. Cuando la cota de volatilidad hace
infactible un perfil, se relaja esa cota manteniendo las dem\'as
restricciones.

\subsection{Frontera eficiente y portafolio de tangencia}

La frontera eficiente se construye resolviendo, para 30 retornos objetivo
$r_{\text{obj}}$ equiespaciados:
\begin{equation}
\min_{w} \quad w^\top \Sigma\, w \qquad \text{s.a.} \quad
\mu^\top w \geq r_{\text{obj}}, \quad
\sum_i w_i = 1, \quad
0 \leq w_i \leq 0.20
\end{equation}

El portafolio de tangencia (``Ideal de mercado'') se selecciona como el punto
de la frontera que maximiza el ratio de Sharpe:
\begin{equation}
\text{Sharpe} = \frac{R_p - R_f}{\sigma_p}
\end{equation}
donde $R_p$ es el retorno anualizado del portafolio y $\sigma_p$ su
volatilidad anualizada. Este portafolio es te\'orico: asume conocimiento
perfecto de los par\'ametros, por lo que su desempe\~no fuera de muestra mide
el grado de \textit{overfitting} de la estimaci\'on hist\'orica.

\subsubsection{Rebalanceo semestral}

El per\'iodo futuro de evaluaci\'on, de 2 a\~nos o aproximadamente 504 d\'ias h\'abiles, se divide en segmentos de 126 d\'ias. Al inicio de cada segmento se actualiza la informaci\'on disponible: se incorporan los retornos ya observados hasta ese punto, se reestiman $\mu$ y $\Sigma$, y se resuelven nuevamente los problemas de optimizaci\'on para todos los perfiles. Los pesos resultantes se aplican al segmento siguiente y se mantienen constantes hasta el pr\'oximo rebalanceo.

Con este procedimiento, el modelo simula una implementaci\'on m\'as realista que una cartera est\'atica, pero sin usar informaci\'on futura no observada. Las m\'etricas finales se calculan concatenando los retornos diarios generados en todos los segmentos de rebalanceo.


% ============================================================================
% 4.4 Modelo Black-Litterman
% ============================================================================
\section{Modelo Black-Litterman}

\subsection{Prior de equilibrio de mercado}

El modelo parte de que, en equilibrio, los retornos impl\'icitos del mercado
son consistentes con los pesos observados por capitalizaci\'on
\cite{black1992, he1999}:
\begin{equation}
\pi = \delta \cdot \Sigma \cdot w_{\text{market}}
\end{equation}
donde $w_{\text{market}}$ son los pesos normalizados por capitalizaci\'on
burs\'atil y $\delta$ es el coeficiente de aversi\'on al riesgo impl\'icito:
\begin{equation}
\delta = \frac{R_{\text{market}} - R_f}{\sigma_{\text{market}}^2}
\end{equation}

\subsection{Incorporaci\'on de \textit{views}}

Dadas las \textit{views} representadas por la matriz de selecci\'on
$P \in \mathbb{R}^{k \times n}$, el vector de expectativas
$Q \in \mathbb{R}^k$ y la matriz de incertidumbre
$\Omega \in \mathbb{R}^{k \times k}$, el retorno esperado posterior es:
\begin{equation}
\mu_{\text{BL}} = \pi + \tau\,\Sigma\, P^\top
\left(P\,\tau\,\Sigma\, P^\top + \Omega\right)^{-1}
(Q - P\,\pi)
\end{equation}
con $\tau = 0.05$. La covarianza posterior se calcula como:
\begin{equation}
\Sigma_{\text{BL}} = (1+\tau)\,\Sigma
- \tau^2\,\Sigma\, P^\top
\left(P\,\tau\,\Sigma\, P^\top + \Omega\right)^{-1}
P\,\Sigma
\end{equation}
y se fuerza a ser semidefinida positiva mediante \texttt{nearest\_psd}. Con
$\mu_{\text{BL}}$ y $\Sigma_{\text{BL}}$ se resuelven los mismos problemas de
optimizaci\'on que en Markowitz.

\subsection{Construcci\'on de \textit{views}}

Se implementan cuatro \textit{views} independientes, cada una generando su
propia versi\'on del modelo Black-Litterman. La
Tabla~\ref{tab:views_config} resume su configuraci\'on.

\begin{table*}[h]
\centering
\caption{Configuración de las cuatro variantes de \textit{views} Black-Litterman.}
\label{tab:views_config}
\scriptsize
\setlength{\tabcolsep}{4pt}
\renewcommand{\arraystretch}{0.9}
\resizebox{\textwidth}{!}{
\begin{tabular}{@{}llccccl@{}}
\toprule
View & Tipo & Lookback & Long & Short & Confianza & Universo \\
\midrule
Momentum & Endógena & 252 días & 20 & 20 & 0.50 & 601 activos \\
Momentum Top20 6M & Market-cap & 126 días & 10 & 10 & 0.50 & Top 20 \\
Momentum Top20-Bottom20 1Y & Market-cap & 252 días & 20 & 20 & 0.50 & Top 40 \\
Desempleo & Macro asumida & 1\,764 días & 466 & 135 & 0.35 & 601 activos \\
\bottomrule
\end{tabular}}
\end{table*}


\subsubsection{Views de momentum}

Las \textit{views} de momentum seleccionan activos por retorno acumulado
durante el per\'iodo de lookback. La matriz $P$ tiene una fila que asigna
$+1/n_{\text{long}}$ a los activos con mayor retorno y
$-1/n_{\text{short}}$ a los de menor retorno. El valor $Q$ corresponde a la
diferencia promedio diaria entre los retornos de ambos grupos. La variante
general opera sobre los 601 activos; las variantes por capitalizaci\'on
restringen el universo a los activos con mayor market cap (top 20 o top 40).

\subsubsection{View de desempleo}

Es un supuesto macroecon\'omico fijo, no un dato observado. Se asume una
tasa de desempleo de 4\% contra un nivel neutral de 5\%. Con desempleo bajo,
se favorecen sectores c\'iclicos (Technology, Consumer Cyclical, Industrials,
Financial Services, Basic Materials, Real Estate, Energy) sobre defensivos
(Healthcare, Consumer Defensive, Utilities, Communication Services). La
se\~nal es proporcional a la diferencia entre ambas tasas, escalada por
$\beta_{\text{macro}} = 1.0$ y convertida a retorno diario dividiendo por
252. La confianza se fija en $c = 0.35$, menor que la de momentum, reflejando
que es un supuesto y no una observaci\'on.

\subsubsection{Calibraci\'on de $\Omega$}

Para cada \textit{view}, la incertidumbre se calibra como:
\begin{equation}
\Omega = \frac{P \cdot (\tau\,\Sigma) \cdot P^\top}{c}
\end{equation}
donde $c$ es la confianza de la \textit{view}. Con una sola \textit{view}
($k=1$), $\Omega$ es escalar y la formulaci\'on se simplifica.

% ============================================================================
% 4.5 Simulación Monte Carlo (P4)
% ============================================================================
\section{Simulaci\'on Monte Carlo (P4)}

\subsection{Objetivo}

La tercera capa traduce los KPIs financieros de Markowitz y Black-Litterman en
proyecciones de riqueza para el cliente y utilidad para la empresa. Simula el
comportamiento completo del sistema recomendador durante 260 semanas (5
a\~nos), incorporando las decisiones del cliente mediante funciones
log\'isticas.

\subsection{Funciones de comportamiento del cliente}

La probabilidad de retiro ante p\'erdidas se modela como:
\begin{equation}
P_1(x_1) = \frac{1}{1 + \exp\left(-s\,(x_1 - \hat{x}_1)\right)}
\end{equation}
donde $x_1 = (C_0 - W_t) / C_0$ es la p\'erdida actual respecto al capital
inicial, $\hat{x}_1$ es la tolerancia m\'axima del perfil, y $s = 20$ es la
sensibilidad log\'istica. $P_1$ se activa \'unicamente cuando
$x_1 > \hat{x}_1$; en caso contrario, la probabilidad de retiro es cero.

La probabilidad de aceptaci\'on de una recomendaci\'on es:
\begin{equation}
P_2(x_2) = \frac{1}{1 + \exp\left(-s\,(x_2 - \hat{x}_2)\right)}
\end{equation}
donde $x_2$ es la tasa de retorno ofrecida en la recomendaci\'on y
$\hat{x}_2$ es el umbral de tolerancia al riesgo del perfil.

\subsection{Estructura de la simulaci\'on}

Para cada combinaci\'on de escenario de datos (con/sin pandemia), t\'ecnica
(6 alternativas), perfil de riesgo (5 niveles) y escenario proyectado
(desfavorable, neutro, favorable), se ejecutan $N = 5\,000$ simulaciones
independientes. En total, esto genera $2 \times 6 \times 5 \times 3 = 180$
combinaciones, cada una con 5\,000 trayectorias.

Cada simulaci\'on opera semana a semana. Primero, el cliente decide si acepta la recomendaci\'on inicial mediante $P_2$; si la acepta, invierte $C_0 = 1\,000$ USD y paga una comisi\'on inicial de $k = 1\%$. Luego, en cada semana se simula el retorno del portafolio como $r_{\text{sem}} \sim \mathcal{N}(\mu_{\text{sem}}, \sigma_{\text{sem}}^2)$, calibrado a partir de los KPIs anualizados de los modelos anteriores. Tambi\'en se ofrece una nueva recomendaci\'on semanal: si el cliente la acepta, nuevamente seg\'un $P_2$, se cobra una comisi\'on de $1\%$ sobre el 5\% de la riqueza actual, que representa la fracci\'on de rotaci\'on asumida. Finalmente, si la p\'erdida acumulada respecto de $C_0$ supera la tolerancia definida para el perfil, se eval\'ua el retiro mediante $P_1$. Cuando el cliente retira, la trayectoria se detiene y la riqueza terminal corresponde al valor disponible en ese momento.


\subsection{Escenarios proyectados}

Los tres escenarios se calibran a partir del retorno anualizado base
$R_{\text{base}}$ y la volatilidad anualizada $\sigma_{\text{base}}$ de cada
t\'ecnica:
\begin{align}
R_{\text{favorable}} &= R_{\text{base}} + 0.5\,\sigma_{\text{base}} \\
R_{\text{neutro}} &= R_{\text{base}} \\
R_{\text{desfavorable}} &= R_{\text{base}} - 0.5\,\sigma_{\text{base}}
\end{align}
Esto genera un abanico sim\'etrico de $\pm 0.5\sigma$ alrededor del caso
base.

\subsection{Reproducibilidad}

La simulaci\'on utiliza semillas deterministas generadas mediante SHA-256 a
partir de la combinaci\'on \'unica de escenario, t\'ecnica, \textit{view} y
perfil. Esto garantiza que cada celda del experimento es reproducible
individualmente sin afectar a las dem\'as.

% ============================================================================
% 4.6 Análisis distribucional de retornos
% ============================================================================
\section{An\'alisis distribucional de retornos}

Complementariamente a la optimizaci\'on, se realiz\'o un an\'alisis
distribucional sobre los retornos diarios del universo F5 sin pandemia, con el
objetivo de caracterizar la forma estad\'istica de los retornos y evaluar la
validez del supuesto de normalidad.

Se ajustaron siete familias param\'etricas candidatas: normal, t-Student,
Laplace, log\'istica, skew-normal, normal generalizada y Johnson~SU. Los
par\'ametros se estimaron por m\'axima verosimilitud. Para cada familia $m$ se
calcularon los criterios de informaci\'on:
\begin{equation}
\text{AIC}_m = 2k_m - 2\ell_m, \qquad
\text{BIC}_m = k_m \log(n) - 2\ell_m,
\end{equation}
donde $k_m$ es el n\'umero de par\'ametros, $n$ la cantidad de retornos y
$\ell_m$ la log-verosimilitud. Se aplicaron adicionalmente pruebas de
normalidad (Shapiro-Wilk, D'Agostino-Pearson, Jarque-Bera) y una prueba
Kolmogorov--Smirnov aproximada.

% ============================================================================
% 4.7 Validación robusta
% ============================================================================
\section{Validaci\'on robusta}

Para evitar que las conclusiones dependan de un \'unico horizonte de entrenamiento, se aplicaron dos mecanismos de validaci\'on. Primero, se evaluaron 70 particiones aleatorias train/test, lo que permiti\'o observar la distribuci\'on del Sharpe fuera de muestra para cada portafolio y t\'ecnica. Segundo, se probaron 18 ventanas m\'oviles robustas, desde h2\_f7 hasta h7\_f2, variando la proporci\'on entre datos de entrenamiento y evaluaci\'on para medir c\'omo cambia el desempe\~no futuro al usar m\'as o menos historia.

El horizonte central seleccionado fue h7\_f2, equivalente a 7 a\~nos de entrenamiento y 2 a\~nos de evaluaci\'on. Este horizonte entrega el mejor Sharpe futuro promedio: 2,32 en Markowitz y resultados comparables en Black-Litterman, con 2,32 para la view de desempleo y 2,22 para momentum. La validaci\'on muestra que ventanas hist\'oricas m\'as largas producen estimaciones m\'as estables, mientras que reducir el periodo de entrenamiento deteriora progresivamente el Sharpe fuera de muestra.

% ============================================================================
% 4.8 Calibración secuencial de Black-Litterman (Entrega 3)
% ============================================================================
\section{Calibraci\'on secuencial de Black-Litterman}

\subsection{Dise\~no de tres horizontes}

La calibraci\'on de los hiperpar\'ametros de Black-Litterman se realiza bajo
un dise\~no experimental que separa estrictamente tres horizontes temporales
no intercambiables. Sea $\mathcal{D}$ el conjunto completo de observaciones
hist\'oricas, compuesto por retornos diarios de $N = 601$ activos. El
experimento divide $\mathcal{D}$ en tres bloques funcionales:

\begin{equation}
\mathcal{D} = \mathcal{H}_1 \cup \mathcal{H}_2 \cup \mathcal{H}_3,
\end{equation}

donde $\mathcal{H}_1$ corresponde al horizonte de calibraci\'on,
$\mathcal{H}_2$ al horizonte de validaci\'on y $\mathcal{H}_3$ al horizonte de
test limpio para la evaluaci\'on econ\'omica P4. La condici\'on fundamental es
que la selecci\'on de hiperpar\'ametros no puede depender de $\mathcal{H}_3$:

\begin{equation}
\widehat{\theta}_{BL} = \mathcal{A}(\mathcal{H}_1), \qquad
\widehat{\theta}_{BL} \not\equiv \mathcal{A}(\mathcal{H}_3).
\end{equation}

Esta separaci\'on evita el problema de \textit{data snooping}: que los
par\'ametros se ajusten al mismo per\'iodo que despu\'es se usa para evaluar
el desempe\~no econ\'omico.

\subsection{Etapas de calibraci\'on}

La calibraci\'on se organiza en cinco etapas secuenciales, cada una evaluada
dentro de $\mathcal{H}_1$ y validada en $\mathcal{H}_2$ antes de congelar la
decisi\'on y avanzar. Las etapas son:

\begin{enumerate}
    \item \textbf{Stage 1 --- Familia de \textit{view}:} Se prueban cuatro
    familias (desempleo macro, momentum general, momentum top/bottom 1Y,
    momentum top market-cap 6M). Para cada una se ejecuta Black-Litterman con
    configuraci\'on base y se selecciona la familia con mayor score de
    calibraci\'on. El ganador en $\mathcal{H}_1$ fue desempleo macro, con
    Sharpe medio 0.854 y delta Sharpe 0.159.

    \item \textbf{Stage 2 --- Estructura de $P$:} Sobre la familia ganadora,
    se prueban combinaciones de par\'ametros de la matriz de selecci\'on $P$:
    tasa de desempleo asumida ($U \in \{3\%, 4\%, 5\%\}$) y sensibilidad
    sectorial ($\beta \in \{1.0, 1.5, 2.0\}$). El ganador fue
    $U = 4\%$ con $\beta = 1.5$.

    \item \textbf{Stage 3 --- Intensidad de $Q$:} Se escala el vector de
    expectativas $Q$ mediante un factor $q_{\text{scale}} \in \{0.5, 1.0, 1.5,
    2.0\}$. El valor \'optimo fue $q_{\text{scale}} = 1.5$.

    \item \textbf{Stage 4 --- Confianza $\Omega$:} Se prueba la confianza
    $c \in \{0.25, 0.35, 0.50, 0.75, 1.00\}$ en la calibraci\'on de $\Omega$.
    Con mayor $c$, menor es la incertidumbre asignada a la \textit{view} y
    mayor su influencia en $\mu_{\text{BL}}$. El ganador fue $c = 0.50$.

    \item \textbf{Stage 5 --- Par\'ametro $\tau$:} Se eval\'ua
    $\tau \in \{0.05, 0.10, 0.20, 0.50, 1.00\}$. El valor $\tau = 0.20$
    result\'o \'optimo.
\end{enumerate}

La configuraci\'on congelada resultante es: \textit{view} de desempleo macro,
$U = 4\%$, $\beta = 1.5$, $q_{\text{scale}} = 1.5$, $c = 0.50$, $\tau = 0.20$.

\subsection{Mecanismo de puntuaci\'on}

En cada etapa, cada configuraci\'on candidata se eval\'ua mediante un score
compuesto que combina el Sharpe medio fuera de muestra en $\mathcal{H}_2$, el
delta de Sharpe respecto a un baseline sin calibrar, y un t\'ermino de
penalizaci\'on por inestabilidad. La configuraci\'on con mayor score acumulado
avanza a la etapa siguiente, y las decisiones previas no se reabren.

% ============================================================================
% 4.9 Validación rolling de tres horizontes (Entrega 3)
% ============================================================================
\section{Validaci\'on rolling de tres horizontes}

La arquitectura de validaci\'on se implementa mediante ventanas m\'oviles
(\textit{rolling windows}) que recorren el tiempo, generando m\'ultiples
particiones independientes de $\mathcal{H}_1$, $\mathcal{H}_2$ y
$\mathcal{H}_3$. Para el escenario con pandemia se usaron 10 ventanas de
calibraci\'on, 4 de test P4 y 3 de validaci\'on; para el escenario sin
pandemia, 3, 1 y 1 respectivamente.

Cada ventana produce una evaluaci\'on completa del pipeline: calibraci\'on de
BL sobre $\mathcal{H}_1$, generaci\'on de portafolios en $\mathcal{H}_2$, y
simulaci\'on Monte Carlo P4 en $\mathcal{H}_3$. Las m\'etricas agregadas
(mejora de Sharpe, delta de drawdown, porcentaje de ventanas donde BL es
recomendado sobre Markowitz base) se reportan por perfil de riesgo y escenario
de datos.

El protocolo tambi\'en registra m\'etricas de din\'amica del portafolio:
turnover medio, n\'umero efectivo de activos ($N_{\text{eff}}$) y concentraci\'on
sectorial (HHI) durante el horizonte de test limpio. Estas m\'etricas
permiten evaluar si la mejora financiera de BL es operacionalmente sostenible
o viene acompa\~nada de rotaci\'on y concentraci\'on excesivas.

% ============================================================================
% 4.10 Simulación de sensibilidad de K% (Entrega 3)
% ============================================================================
\section{Simulaci\'on de sensibilidad de la comisi\'on $K\%$}

\subsection{Objetivo}

El objetivo es determinar el valor de la comisi\'on $K\%$ que maximiza el
Score P4 del recomendador, entendiendo c\'omo afecta simult\'aneamente la
riqueza del cliente, la utilidad de FinPUC y la tasa de retiro.

\subsection{Dise\~no experimental}

Se eval\'uan 8 valores de comisi\'on: $K \in \{0.25\%, 0.5\%, 0.75\%,
1\%, 1.5\%, 2\%, 3\%, 5\%\}$. Para cada $K$, se ejecuta la simulaci\'on Monte
Carlo completa: 5\,000 trayectorias $\times$ 260 semanas $\times$ 5 perfiles
$\times$ 4 metodolog\'ias (Equiponderado, Markowitz, M\'inima Varianza, BL
Momentum General calibrado). Los retornos semanales se simulan como normales
i.i.d.\ con media y volatilidad provenientes de los KPIs de la Entrega 2,
usando el escenario proyectado neutro (sin shift $\pm 0.5\sigma$).

La comisi\'on se aplica en dos momentos: (a) comisi\'on inicial de $K\%$ sobre
$C_0 = 1\,000$ USD al aceptar la primera recomendaci\'on, y (b) comisi\'on de
$K\%$ sobre el 5\% de la riqueza actual en cada recomendaci\'on aceptada
posteriormente, donde el 5\% representa la fracci\'on de rotaci\'on asumida.

\subsection{M\'etricas}

Para cada combinaci\'on $(K, \text{metodolog\'ia}, \text{perfil})$ se calculan:
riqueza terminal media, tasa de retiro media, utilidad empresa media, y el
Score P4 definido en la Ecuaci\'on~(\ref{eq:score_p4}). El $K$ \'optimo es
aquel que maximiza el Score P4 promedio sobre todas las metodolog\'ias y
perfiles.

% ============================================================================
% 4.11 Optimización multiobjetivo de λ (Entrega 3)
% ============================================================================
\section{Optimizaci\'on multiobjetivo de $\lambda$}

\subsection{Objetivo}

Determinar el par\'ametro de balance $\lambda \in [0,1]$ que optimiza el
desempe\~no conjunto del recomendador, ponderando expl\'icitamente el retorno
del cliente y la utilidad de la empresa.

\subsection{Formulaci\'on}

Para cada $\lambda$, se define el score compuesto:
\begin{equation}
\label{eq:lambda_score}
S_\lambda(w) = \lambda \cdot f_1(w) + (1-\lambda) \cdot f_2(w),
\end{equation}
donde $f_1(w) = \mathbb{E}[w^\top r]$ es el retorno esperado anualizado del
portafolio y $f_2(w) \propto K \cdot P_2(\text{return}) \cdot
(1 - P_1(\text{loss}))$ es la utilidad esperada de FinPUC. Ambas funciones se
normalizan mediante \textit{min-max scaling} al rango $[0,1]$ antes de
combinar, usando los valores m\'inimos y m\'aximos observados en todo el
experimento.

\subsection{Dise\~no experimental}

Se eval\'uan 11 valores de $\lambda$ en grilla discreta:
$\lambda \in \{0, 0.1, 0.2, 0.3, 0.4, 0.5, 0.6, 0.7, 0.8, 0.9, 1.0\}$.
Cada $\lambda$ se prueba con el $K$ \'optimo determinado en la secci\'on
anterior y con la metodolog\'ia de mejor Score P4. Para cada $\lambda$, se
ejecuta Monte Carlo con 5\,000 trayectorias y se calcula el Score P4
resultante, la riqueza terminal y la utilidad empresa.

Los resultados se presentan como una frontera de Pareto cliente vs.\ empresa,
donde cada $\lambda$ define un punto en el espacio $\big(f_1(w),
f_2(w)\big)$. El $\lambda$ \'optimo es aquel que maximiza el Score P4
observado, y se compara con el $\lambda$ que maximiza un score ponderado que
asigna igual importancia a riqueza y utilidad normalizadas.

% ============================================================================
% 4.12 Comparación de modelos de abandono y views (Entrega 3)
% ============================================================================
\section{Comparaci\'on de modelos de abandono y \textit{views}}

\subsection{Modelos de probabilidad de abandono}

Se comparan tres versiones del modelo de abandono:

\begin{itemize}
    \item \textbf{v1:} Modelo base con activaci\'on de $P_1$ por tolerancia de
    p\'erdida del perfil y $P_2$ por umbral de retorno ofrecido, con
    sensibilidad $s = 20$.

    \item \textbf{v1$+$:} Extiende v1 incorporando una restricci\'on de ruina
    homog\'enea para todas las simulaciones: si la riqueza del cliente cae por
    debajo del 10\% de $C_0$, el abandono es forzoso. Esto homogeniza el
    tratamiento de escenarios extremos entre \textit{views} y permite aislar
    el efecto de la calibraci\'on de la \textit{view} de diferencias en la
    regla de abandono.

    \item \textbf{v2:} Reemplaza el gatillo de p\'erdida respecto a $C_0$ por
    \textit{drawdown} respecto al m\'aximo hist\'orico de riqueza,
    representando mejor la percepci\'on de p\'erdida del cliente.
\end{itemize}

Los tres modelos se eval\'uan bajo el mismo protocolo de validaci\'on de tres
horizontes, y se comparan en t\'erminos de tasa de retiro, riqueza terminal y
estabilidad del ranking P4 entre \textit{views}.

\subsection{Comparaci\'on de \textit{views}}

Las cuatro familias de \textit{views} se comparan bajo el modelo de abandono
v1$+$ y la configuraci\'on BL calibrada:

\begin{itemize}
    \item \textbf{Desempleo macro:} tasa de desempleo asumida de 4\% frente a
    5\% neutral, con $\beta = 1.5$. Favorece sectores c\'iclicos sobre
    defensivos.

    \item \textbf{Momentum general:} 20 activos long y 20 short sobre los 601
    activos, con \textit{lookback} de 252 d\'ias.

    \item \textbf{Momentum top/bottom 1Y:} 20 long y 20 short sobre el top 40
    por capitalizaci\'on, con \textit{lookback} de 252 d\'ias.

    \item \textbf{Momentum top market-cap 6M:} 10 long y 10 short sobre el top
    20 por capitalizaci\'on, con \textit{lookback} de 126 d\'ias.
\end{itemize}

La comparaci\'on se concentra en tres dimensiones: (a) mejora de Sharpe fuera
de muestra, (b) mejora de Score P4 frente a Markowitz base, y (c) estabilidad
operacional medida por turnover, concentraci\'on sectorial (HHI) y n\'umero
efectivo de activos ($N_{\text{eff}}$).
